\section{Question 4: éléments matériels de la contrefaçon}

A3 présente un instrument à commande manuelle pour seringue comportant
un piston pourvu d'un siège adapté pour recevoir un axe de commande du
coulissement du piston à l'intérieur du corps de la seringue
("positionnement (et le retrait) sans difficulté de la tête de la tige
filetée à l'intérieur du logement d'un piston de seringue" par.2),
l'instrument comprenant:
\begin{itemize}
    \item un moyen d'accouplement au siège du piston ("une tête conique
    " par.2);
    \item un axe fileté portant un filetage extérieur (figure);
    \item un écrou d'actionnement manuel pourvu d'un trou traversant
    portant un filetage intérieur ayant le même pas que le filetage
    de l'axe fileté et au travers duquel l'axe fileté passe et vient
    en prise (figure)
    \item une cinématique de sorte que lorsque le moyen d'accouplement
    est couplé au siège du piston, la rotation de l'écrou d'actionnement
    par rapport au corps de la seringue provoque le coulissement axial
    de l'axe fileté au travers de son trou traversant et le déplacement
    du piston à l'intérieur du corps de la seringue ("actionner manuellement
    la rotation de la viole en sens anti-horaire de façon à entraîner
    un déplacement axial de la tige filetée et du piston de la seringue" pt.4
    du mode opératoire).\\
\end{itemize}

A3 ne présente pas:
\begin{itemize}
    \item un axe fileté en extrémité duquel le moyen d'accouplement
    est solidement fixé.\\
\end{itemize}

En effet, \epcart{69}(1) CBE dispose que les revendications sont interprétées en fonction de la description. Ainsi,
le terme "solidement fixé" s'interprète à la lumière de la description:
"\textit{De préférence, les moyens d'accouplement 13 de l'instrument 1 au siège 7 du piston 5 comprennent
une tête filetée 13a pour cette fixation qui soit suffisamment solide compte tenu, d'une part,
de l'effort de traction exercé par l'axe 9 qui contrôle le déplacement du piston 5 et, d'autre part,
de la résistance exercée à l'encontre du coulissement de ce piston 5.}" (par.[0011]). Or,
dans A3, c'est précisément l'inverse qui est décrit, un effort de traction détache
l'instrument du siège du piston (figure 7 du mode opératoire).
Donc, A3 ne divulgue pas un axe fileté
en extrémité duquel le moyen d'accouplement est \underline{solidement fixé}.\\

A3 présente l'instrument défini à la revendication 1, l'instrument comprenant
une virole annulaire ("un disque de fixation" par.3):
\begin{itemize}
    \item un trou traversant au travers duquel coulisse l'axe fileté (figure);
    \item une surface d'engagement avec le bord de l'ouverture d'extrémité
    du corps de la seringue ("la fixation peut s'opérer par clipsage,
    vissage ou baïonnette. L'assemblage du dispositif avec le réservoir de
    la seringue, par l'intermédiaire du disque de fixation" par.3);
    \item une surface de forme complémentaire à une surface de butée tronconique
    de l'écrou d'actionnement ("le disque de fixation est monté fou sur l'écrou"
    par.3).\\
\end{itemize}

A3 présente un procédé de mise en œuvre d'un instrument à commande manuelle
avec une seringue comprenant un piston pourvu d'un siège, le procédé comprenant:
\begin{itemize}
    \item une étape de couplage de l'instrument à la seringue par le couplage
    du moyen d'accouplement au siège du piston ("positionner la tête du dispositif
    dans le logement du piston de la seringue" pt.1 du mode opératoire), complétée
    par un couplage de la viole annulaire avec l'ouverture du corps de la seringue
    ("verrouiller le disque de fixation à l'extrémité du réservoir de la seringue"
    pt.3 du mode opératoire);
    \item un déplacement de l'écrou d'actionnement en butée contre la viole annulaire
    positionnée entre l'écrou d'actionnement et le corps de la seringue
    ("ajuster la position du disque de fixation par pivotement en sens horaire
    de la viole le long de la tige filetée" pt.2 du mode opératoire).
    \item une rotation de l'écrou d'actionnement en butée par rapport au corps
    de la seringue, provoquant ainsi le coulissement axial de l'axe fileté à
    l'intérieur du trou de l'écrou d'actionnement.\\
\end{itemize}

A3 ne présente pas:
\begin{itemize}
   \item une traction du piston vers
    l'ouverture d'extrémité du corps pour réaliser une aspiration au travers de la
    pointe de la seringue.\\
\end{itemize}

\textbf{Conclusion: le dispositif WHITE ne reproduit pas les revendications
de la partie française du brevet européen EP01.}